\subsection{Excluding the trivial residual dyadic comparison}

\begin{lemma}[B-RD46: the cubic part of dyadic elliptic inertia]
\label{lem:B-RD46-discriminant}
\status{Proved; complete independent parent and obstructions/crystals proof and exact-code audits passed.}
Let $K$ be a complete discretely valued field of characteristic zero
with perfect residue field of characteristic two. Suppose that an
elliptic curve $E/K$ is potentially supersingular. If $\delta$ is the
valuation of the discriminant of any Weierstrass equation over $K$,
then the order of the finite special-inertia image is divisible by
three if and only if $3\nmid\delta$. Consequently its character on
elements of odd order is determined by whether $3\mid\delta$.
Here ``odd order'' refers to the image in a common finite
special-inertia quotient.
\end{lemma}
\begin{proof}
Pass to the completed maximal unramified extension, which preserves
the valuation and the special-inertia representation. Let $L/K$ be
the minimal Galois extension over which the curve has good reduction,
with faithful special action $H$ and $m=[L:K]$. The explicit
supersingular automorphism group and finite-subgroup classification
in Lemma~\ref{lem:B-RD44-conjugacy} show that
$m\in\{1,2,3,4,6,8,24\}$.
If $\omega$ is the invariant differential of the equation over $K$
and $\omega_L$ a good-model invariant differential, write
$\omega=c\omega_L$. For a coordinate scaling $x=u^2x_0$, $y=u^3y_0$, one has
$\Delta=u^{12}\Delta_0$ and $\omega=u^{-1}\omega_0$.
Thus the discriminant transformation rule gives
\[
                         12v_L(c)=-m\delta.             \tag{B-RD46.1}
\]
This identity is valid even if the original equation is not integral.
If $3\nmid m$, it immediately gives $3\mid\delta$.

Suppose $3\mid m$ and choose $\tau\in H$ of order three.
Its action on a uniformizer satisfies
$\overline{\tau(\pi_L)/\pi_L}=\zeta$, a primitive cube root:
the kernel of the tame cotangent character is a two-group.
Since the residue action is trivial, the multiplier of $\tau$ on
the good special differential is $\zeta^{-v_L(c)}$.
This multiplier is nontrivial. Indeed an automorphism of order
three with derivative one at the origin can be linearized by
averaging its local parameter, using that three is invertible;
it would then be the identity on the formal completion and hence
on the elliptic curve. Faithfulness excludes that possibility.
Thus $3\nmid v_L(c)$, and (B-RD46.1), with $v_3(m)=1$, gives
$3\nmid\delta$.

The only odd orders in $H$ are one and three. An order-three
automorphism has eigenvalues $\zeta_3,\zeta_3^{-1}$ on the
rational special Dieudonne module and trace $-1$; the identity
has trace two. If the cubic quotient is nontrivial, its kernel
on tame inertia is the unique subgroup of index three.
Hence this determines the character as a function on inertia,
not merely its abstract character table.
\end{proof}

\begin{proposition}[B-RD46: rational two-torsion excludes the ambiguity]
\label{prop:B-RD46}
\status{Proved; complete independent parent and obstructions/crystals proof and exact-code audits passed.}
Let $K$ be a complete discretely valued field of mixed characteristic
$(0,2)$, with perfect residue field and $v_K(2)=1$. Let $E,E'/K$ be
potentially supersingular elliptic curves with rational full
two-torsion. If $E[4]\simeq E'[4]$ over $K$, their special
inertia characters agree on elements of odd order.

In particular, in the remaining height-two, ramification-one
sixteen-torsion comparison of Proposition~\ref{prop:B-RD44},
the common residual inertia image cannot be trivial. After
completed maximal unramified base change, it is either $C_2$
or $C_3$. This proposition decides neither remaining case.
\end{proposition}
\begin{proof}
Every curve with rational full two-torsion is a quadratic twist
of a Legendre curve
\[
              E_\lambda:\quad y^2=x(x-1)(x-\lambda),
              \qquad \lambda\in K\setminus\{0,1\}.
\]
Permuting the three nonzero two-torsion points permits the
normalization
\[
 v_K(\lambda)=q\ge0,\qquad v_K(1-\lambda)=0.
                                                               \tag{B-RD46.2}
\]
If $q=0$, both differences are units. If $q>0$, the formula
\[
 j(E_\lambda)=2^8\frac{(1-\lambda+\lambda^2)^3}
                         {\lambda^2(1-\lambda)^2}
\]
gives $v_K(j)=8-2q$. Potential supersingularity requires
$v_K(j)>0$, so $q\in\{1,2,3\}$ in this case.
A twist by $d\in K^\times$ has discriminant
$2^4d^6\lambda^2(1-\lambda)^2$, and therefore
\[
 \delta\equiv4+2q\pmod3.
 \quad
 \begin{array}{c|rrrr}
 q&0&1&2&3\\\hline
 \delta\bmod3&1&0&2&1
 \end{array}                                             \tag{B-RD46.3}
\]
The preceding lemma shows that different odd-order characters
would require one normalized parameter to have $q=1$ and the
other to have $q'=0,2$, or $3$.

We compute the complete projective four-torsion action. Choose
\[
 i^2=-1,\qquad a^2=\lambda,\qquad b^2=1-\lambda,
\]
and take the basis
\[
 P=(a,ia(a-1)),\qquad Q=(1+b,b(1+b)).
                                                               \tag{B-RD46.4}
\]
The chord-and-tangent law gives
$2P=(0,0)$ and $2Q=(1,0)$.
Changing the sign of $i$ negates $P$; changing the sign of $a$
sends $P$ to $P+2Q$; changing the sign of $b$ sends $Q$ to
$Q+2P$. These statements remain valid when the three quadratic
characters are dependent: their simultaneous effects commute
modulo four. Write $\chi_z:G_K\to\mathbf F_2$ for the
quadratic Kummer character of $z$, and set
\[
 N=\begin{pmatrix}0&0\\1&0\end{pmatrix},\quad
 U=\begin{pmatrix}0&1\\0&0\end{pmatrix},\quad
 T=\begin{pmatrix}1&0\\0&0\end{pmatrix}.
\]
For the quadratic twist by $d$ the exact action is
\[
 \rho_4(\sigma)=I+2A(\sigma),\qquad
 A(\sigma)=\chi_d(\sigma)I+\chi_{-1}(\sigma)T
             +\chi_\lambda(\sigma)N
             +\chi_{1-\lambda}(\sigma)U.              \tag{B-RD46.5}
\]
Thus the scalar term accounts for every possible quadratic twist.
An isomorphism of the two four-torsion modules conjugates their
matrices $A$ by one fixed $C\in\GL_2(\mathbf F_2)$.

The valuation modulo two is a linear functional on
$K^\times/K^{\times2}$, zero on unit classes and nonzero on
the class of two. Kummer theory therefore supplies an element
$\sigma_0\in G_K$ fixing all unit square roots needed in both
displays and changing the sign of $\sqrt2$.
It is enough to apply this to the finite-dimensional space of
classes occurring in the two equations and twists; there is no
assumption that the full square-class group is finite.
For a normalized parameter of valuation $q$, evaluation at
$\sigma_0$ gives
\[
                       A(\sigma_0)=\epsilon I+(q\bmod2)N.
                                                               \tag{B-RD46.6}
\]
If $q=1$ and $q'=0$ or $2$, one of these matrices is nonscalar
and the other scalar, so they cannot be conjugate.

It remains to compare $q=1$ and $q'=3$.
Equation (B-RD46.6) forces
$CNC^{-1}=N+(\epsilon+\epsilon')I$.
The determinant distinguishes $N$ from $N+I$, so
$\epsilon=\epsilon'$ and $C$ centralizes $N$.
The centralizer in $\GL_2(\mathbf F_2)$ is exactly
$\{I,I+N\}$. Conjugation by either matrix preserves the
upper-right entry of every matrix. Taking that entry in
(B-RD46.5), for every $\sigma$, now forces
\[
                         \chi_{1-\lambda}
                              =\chi_{1-\lambda'}.       \tag{B-RD46.7}
\]
The right side is zero: since $v_K(\lambda')=3$,
Hensel's lemma applied to $X^2-(1-\lambda')$ at $X=1$
gives a square root, using $3>2v_K(2)=2$.
The left side is nonzero. If $(1+w)^2=1-\lambda$, reduction
forces $v_K(w)\ge1$, whereas
$v_K(2w+w^2)\ge2$ contradicts $v_K(\lambda)=1$.
This contradiction proves the proposition.

For its final assertion, use the exact elliptic realizations
of Proposition~\ref{prop:B-RD44-realization} and restrict the
given sixteen-torsion isomorphism to four-torsion. A trivial
residual inertia action gives rational full two-torsion after
completed maximal unramified base change, which has just been
excluded. Finally the wild image in $\GL_2(\mathbf F_2)=S_3$
is a normal two-subgroup and its tame quotient is cyclic of
odd order. This excludes $S_3$ itself and leaves, among the
nontrivial images, precisely $C_2$ and $C_3$.
\end{proof}

\paragraph{Exact certificate and scope.}
The retained program
\path{code/verify_p2_residual_trivial_gate.py}
checks the displayed halving and translation identities over an
exact polynomial quotient, the six conjugating matrices, their
centralizer and upper-right-entry property, and the valuation
table. It does not enumerate elliptic curves or decide the
remaining residual $C_2$ and $C_3$ cases. Both new assertions have passed complete independent audits.

\paragraph{Self-check correction.}
The initial written draft put the opposite sign in (B-RD46.1).
The coordinate-scaling check above corrected it before independent
audit; the original is retained at
\path{failed_candidates/B46_differential_sign_v1.tex}.
Only nondivisibility by three is used, so this correction changes
no conclusion. The exact certificate passed in
$0.12948673400387634$ seconds.

\paragraph{Independent audits.}
The parent independently checked the entire proof and code, with a
preserved-output rerun of $0.35755470999720274$ seconds recorded
at 2026-11-19 14:16:15.762965 UTC. The obstructions/crystals agent
completed the same full audit, rerunning the certificate in
$0.3844016500006546$ seconds. The prior B44 original-source
verification supplies the finite special-inertia classification;
the remaining new steps are proved explicitly here.

\paragraph{Subsequent closure.}
The remaining residual branches are settled in
Theorem~\ref{thm:B-RD48}, which gives the exact value
$\nu_{\min}(2,1,2)=3$. The assertions and finite certificate
in this subsection retain their stated narrower scope.
